\section{Example prompts and traces}\label{sec:app_examples}

This appendix provides detailed examples of the LLM prompts and intermediate system states generated during the SolEvolve execution cycles. These verbose traces offer deeper insight into the semantic interpretation mechanics underlying the Reflector and Generator agents.

\subsection{Reflector feedback example: stagnation diagnosis}\label{sec:app_stagnation}

The following example illustrates the Reflector's operation when the Evolver's population stagnates at a local optimum (discussed in Section~\ref{ch:design}).

\begin{tcolorbox}[enhanced, breakable, colback=white, colframe=black, title={Reflector Feedback Example: Stagnation Diagnosis}]
\textbf{Input Signal:} $d_{\min} = 12$ stable for 50 generations

\tcbline

\textbf{Semantic Interpretation:}\\
\textit{``Local optimum detected in weight-12 codeword space. Population has converged to structurally similar candidates sharing weight-12 attractors.''}

\tcbline

\textbf{Generated Intervention Prompt:}\\
\textit{``Propose symmetry-breaking constraints to eliminate weight-12 codewords, or suggest alternative encoding that penalizes low-weight solutions. Consider column-swap mutations biased toward positions where weight-12 codewords concentrate.''}

\tcbline

\textbf{State Transition:} Exploitation $\to$ Repair (trigger SAT-based correction)
\end{tcolorbox}


\subsection{Generator prompt: SAT encoding strategy selection}\label{sec:app_generator}

The following prompt template forces the Generator to perform a structured trade-off analysis when selecting a SAT encoding for cardinality constraints.

\begin{tcolorbox}[enhanced, breakable, colback=white, colframe=black, title={Generator Prompt: SAT Encoding Strategy Selection}]
\textbf{Task:} SAT Encoding Strategy for $[n,k,d]$ Code

\tcbline

\textbf{Target Parameters:}
\begin{itemize}
\item Binary linear code: $[\texttt{\{n\}},\texttt{\{k\}},\texttt{\{d\}}]$
\item Variables: $k(n-k) = \texttt{\{num\_vars\}}$ parity bits
\item Non-zero codewords: $2^k - 1 = \texttt{\{num\_codewords\}}$
\item Constraint: Each codeword must have Hamming weight $\geq d$
\end{itemize}

\tcbline

\textbf{Encoding Options} (evaluate each):
\begin{enumerate}
\item \textbf{Binomial:} $O\binom{n}{d-1}$ clauses, arc-consistent propagation
\item \textbf{Sequential counter:} $O(nd)$ clauses, weaker propagation
\item \textbf{Commander:} $O(n\sqrt{n})$ clauses, balanced trade-off
\end{enumerate}

\tcbline

\textbf{Required Output Format:}

\texttt{<analysis>} Compare clause counts, propagation strength, and memory footprint for this specific instance \texttt{</analysis>}

\texttt{<recommendation>} Selected encoding with explicit justification \texttt{</recommendation>}

\texttt{<implementation>} Python function signature and key logic \texttt{</implementation>}
\end{tcolorbox}


\subsection{Detailed analysis: discovery of $[22,11,7]$ code}\label{sec:app_case_study}

This example provides the full computational trace for the hybrid SAT-GA discovery discussed in Section~\ref{sec:hybrid}.

\begin{tcolorbox}[enhanced, breakable, colback=white, colframe=black, title={Detailed Analysis: Discovery of {[22,11,7]} Code}]
\textbf{Discovery Process:}
\begin{enumerate}
\item \textbf{SAT Seeding:} Generated 20 inequivalent binary $[22,11,6]$ codes to ensure diverse initial population
\item \textbf{GA Initialization:} Seeded population of 100 individuals (20 SAT-derived $+$ 80 mutated)
\item \textbf{Evolution:} Pure GA ran for 50 generations $\to$ stagnated at $d_{\max} = 6$
\item \textbf{SAT Repair:} At generation 50, repair module triggered $\to$ eliminated codewords with weight $< 7$
\item \textbf{Result:} Solver successfully repaired one candidate $\to$ valid $[22,11,7]$ code
\end{enumerate}

\tcbline

\textbf{Structural Properties:}\\[0.3em]
\begin{tabular}{@{}ll@{}}
\textit{Weight variance:} & $5.44$ (vs. 21.0 in standard codes) \\
\textit{Quality indicator:} & High structural uniformity
\end{tabular}

\tcbline

\textbf{Generator Matrix} $G = [I_{11} \mid P]$:
\begin{center}
\resizebox{0.9\linewidth}{!}{%
$G = \left[
\begin{array}{c|ccccccccccc}
 & 1 & 0 & 0 & 0 & 1 & 0 & 1 & 1 & 1 & 1 & 1 \\
 & 1 & 0 & 1 & 1 & 1 & 0 & 1 & 1 & 0 & 0 & 0 \\
 & 1 & 0 & 1 & 0 & 0 & 1 & 1 & 0 & 0 & 1 & 1 \\
 & 1 & 1 & 0 & 1 & 0 & 1 & 1 & 1 & 0 & 1 & 0 \\
 & 0 & 0 & 1 & 1 & 1 & 1 & 1 & 0 & 1 & 1 & 0 \\
I_{11} & 0 & 1 & 1 & 0 & 0 & 1 & 0 & 1 & 1 & 1 & 1 \\
 & 0 & 1 & 1 & 0 & 1 & 0 & 1 & 0 & 1 & 0 & 1 \\
 & 1 & 1 & 0 & 0 & 1 & 1 & 0 & 1 & 0 & 0 & 1 \\
 & 1 & 1 & 1 & 1 & 0 & 1 & 0 & 0 & 1 & 0 & 0 \\
 & 0 & 1 & 0 & 1 & 1 & 0 & 0 & 1 & 1 & 1 & 0 \\
 & 1 & 1 & 1 & 1 & 1 & 0 & 0 & 0 & 0 & 1 & 1
\end{array}
\right]$%
}
\end{center}

\tcbline

\textbf{Weight Enumerator Polynomial:}\\
\begin{align*}
W(z) =\; & 1 + 176z^7 + 330z^8 \\ 
&+ 672z^{11} + 616z^{12} + 176z^{15} \\
&+ 77z^{16}.
\end{align*}
\end{tcolorbox}
